% !PRINTED 32
\weberpage{32}
Demnach ist unser ganzes Problem enthalten in der Gleichung (6), die
nichts anderes besagt, als daß die Funktion $4n$ten Grades $f(U,V)$,
$2n-2$ quadratische Faktoren enthalten soll. Zur Befriedigung der
hieraus folgenden $(2n-2)$ Bedingungen hat man die $2n+2$ in $U,V$
enthaltenen Koeffizienten zur Verfügung, so daß vier von diesen
unbestimmt bleiben. Dies war vorauszusehen, da für $\xi,\eta$ beliebige
homogene lineare Funktionen von $\xi,\eta$ eingeführt werden können. Man
kann diese vier überzähligen Konstanten dazu verwenden, um das
Differential (8) in eine Normalform zu bringen. Da aber die
Bedingungsgleichungen für die Koeffizienten von $U,V$ nicht linear sind,
so gibt es für einen gegebenen Transformationsgrad mehrere
Transformationen.

Wir werden diesem Transformationsproblem später von einer ganz anderen
Seite her wieder begegnen und weit tiefer darauf eingehen müssen. Es
soll daher auf die Einzelheiten des algebraischen Problems hier nicht
näher eingegangen werden; dagegen wollen wir durch zwei Beispiele das
Gesagte veranschaulichen.

\begin{center}
\S\,9. \textbf{\gesperrt{Die Transformation zweiten Grades.}}
\end{center}

Wir setzen, um die elliptischen Differentiale in der Normalform zu
erhalten:
\[
\begin{aligned}
f(x,y) &= xy(x-y)(x-\lambda^2 y),\\
\varphi(\xi,\eta) &= \xi\eta(\xi-\eta)(\xi-\varkappa^2\eta),
\end{aligned}
\]
und es seien $U, V$ vom zweiten Grade.

Die Gleichung (6), \S\,8, wird jetzt, da $T$ vom zweiten Grade ist
\weq{1}{UV(U-V)(U-\lambda^2 V) = (a\xi+b\eta)^2(a'\xi+b'\eta)^2\,\xi\eta(\xi-\eta)(\xi-\varkappa^2\eta),}
und es müssen zwei der vier Faktoren zweiten Grades auf der linken Seite
dieser Gleichung Quadrate sein. Die große Zahl der hierin liegenden
Möglichkeiten wollen wir dadurch noch beschränken, daß wir voraussetzen,
$\eta$ und $y$ sollen gleichzeitig verschwinden, also $V$ durch $\eta$
teilbar sein. Dies gibt (von einem konstanten Faktor abgesehen) für $V$
die folgenden drei Möglichkeiten:
\[
1.\ V = \xi\eta, \qquad 2.\ V = \eta(\xi-\eta), \qquad 3.\ V = \eta(\xi-\varkappa^2\eta).
\]
Jeder dieser drei Fälle umfaßt nun wieder drei Unterfälle, indem von den
drei übrigen Faktoren $U$, $U-V$, $U-\lambda^2 V$ irgend zwei als
Quadrate angenommen werden können. Von diesen letzteren
%CONT>>
